
Выведем полезное следствие, дающее достаточное условие сходимости с вероятностью 1.

\mkcor{Достаточное условие сходимости почти наверное}{suff_as}{
    Пусть $\{\xi_n, n \in \mathbb{N}\}$ --- последовательность с.в., причем для любого $\varepsilon > 0$ сходится ряд
    $$ \sum_{n=1}^{+\infty} \P(|\xi_n - \xi| > \varepsilon) < +\infty. $$
    Тогда $\xi_n \toAS \xi$.
}

\mkproof{
    Проверим критерий сходимости с вероятностью 1:
    $$ \P\left(\sup_{k \ge n}|\xi_k - \xi| > \varepsilon\right) = \P\left(\bigcup_{k \ge n}\{|\xi_k - \xi| > \varepsilon\}\right) \le \sum_{k=n}^{+\infty} \P(|\xi_k - \xi| > \varepsilon) \longrightarrow 0 \quad \text{при } n \to \infty $$
    как остаток сходящегося ряда.
}

Обсуждение предельных теорем мы начнем с закона больших чисел.

\mkthm{Закон больших чисел в форме Чебышёва}{lln_cheb}{
    Пусть $\{\xi_n, n \in \mathbb{N}\}$ --- последовательность попарно некоррелированных с.в., причем дисперсии их равномерно ограничены: для некоторой $C > 0$ выполнено $\forall k \; \D\xi_k \le C$. Обозначим $S_n = \xi_1 + \ldots + \xi_n$. Тогда для любого $\delta > 0$
    $$ \frac{S_n - \E S_n}{n^{1/2+\delta}} \xrightarrow{\P} 0. $$
}

\mkproof{
    Оценим $\D S_n$. В силу попарной некоррелированности
    $$ \D S_n = \D(\xi_1 + \ldots + \xi_n) = \sum_{i=1}^n \D\xi_i \le Cn. $$
    Из неравенства Чебышёва получаем: для любого $\varepsilon > 0$
    $$ \P\left(\left| \frac{S_n - \E S_n}{n^{1/2+\delta}} \right| \ge \varepsilon\right) = \P\left(|S_n - \E S_n| \ge \varepsilon n^{1/2+\delta}\right) \le \frac{\D S_n}{\varepsilon^2 n^{1+2\delta}} \le \frac{C}{\varepsilon^2 n^{2\delta}} \longrightarrow 0. $$
}

Сформулируем простое следствие, несколько проясняющее смысл закона больших чисел.

\mkcor{}{cor_lln_cheb}{
    Пусть $\{\xi_n, n \in \mathbb{N}\}$ --- независимые с.в., имеющие одинаковые средние $\E\xi_k = a$ и равномерно ограниченные дисперсии: для некоторой $C > 0$ выполнено $\forall k \; \D\xi_k \le C$. Тогда
    $$ \frac{\xi_1 + \ldots + \xi_n}{n} \xrightarrow{\P} a \quad \text{при } n \to \infty. $$
}

\mkproof{
    Достаточно положить $\delta = 1/2$ и заметить, что $\E S_n = na$.
}

\begin{remark}[Смысл закона больших чисел и устойчивость частот]\label{rem:slln_meaning}
    Следствие проясняет смысл закона больших чисел, как первую попытку теоретического обоснования принципа устойчивости частот. Пусть $\xi_1, \ldots, \xi_n, \ldots$ --- независимые наблюдения одного и того же эксперимента. Тогда по ЗБЧ их среднее арифметическое сходится к среднему значению одного результата $\E\xi_i = a$. Если, например, $\xi_i$ --- это индикаторы появления некоторого события $A$ в эксперименте, т.е.
    $$ \xi_i = \I\{A \text{ произошло в } i\text{-ом эксперименте}\}, $$
    то $\nu_n(A) = \frac{\xi_1 + \ldots + \xi_n}{n}$ --- это в точности частота появления $A$, а $\E\xi_n = \P(A)$ --- это вероятность события $A$. Тогда ЗБЧ гласит, что
    $$ \nu_n(A) \xrightarrow{\P} \P(A). $$
    Это очень похоже на феномен устойчивости частот, который постулировался нами в начале как отправная точка изучения теории вероятностей. Правда, закон больших чисел говорит о сходимости по вероятности, что не совсем отвечает естественной интуиции при обсуждении принципа устойчивости частот.

    Естественно было бы видеть утверждение в духе: бросаем монетку, считаем частоту выпадения решки, она стремится к $1/2$, т.е. просто стремится, а не вероятность получить частоту, отличающуюся от $1/2$ хотя на $0.01$, стремится к нулю. С математической точки зрения мы хотели бы заменить сходимость по вероятности в ЗБЧ на более сильную сходимость с вероятностью 1, почти наверное.
\end{remark}

Этим мы и займемся в теореме об усиленном законе больших чисел.

\mkthm{УЗБЧ для некоррелированных с.в.}{slln_uncorr}{
    Пусть $\{\xi_n, n \in \mathbb{N}\}$ --- последовательность попарно некоррелированных с.в., причем их дисперсии равномерно ограничены: для некоторой $C > 0$ выполнено $\forall k \; \D\xi_k \le C$. Обозначим $S_n = \xi_1 + \ldots + \xi_n$. Тогда
    $$ \frac{S_n - \E S_n}{n} \toAS 0. $$
}

\mkproof{
    Без ограничения общности, считаем, что $\E\xi_n = 0$ для всех $n \in \mathbb{N}$. Покажем сначала, что $S_n/n$ сходится почти наверное к нулю по подпоследовательности $n^2$. Действительно, для любого $\varepsilon > 0$
    $$ \P\left(\left| \frac{S_{n^2}}{n^2} \right| \ge \varepsilon\right) \le |\text{неравенство Чебышева}| \le \frac{\D S_{n^2}}{\varepsilon^2 n^4} = |\text{некоррелированность}| = \frac{\sum_{k=1}^{n^2} \D\xi_k}{\varepsilon^2 n^4} \le \frac{C}{\varepsilon^2 n^2}. $$
    Значит, ряд $\sum_{n=1}^\infty \P(|S_{n^2}/n^2| \ge \varepsilon)$ сходится и, следовательно, $S_{n^2}/n^2 \toAS 0$.

    Для $n \in \mathbb{N}$ определим $m(n) = \lfloor \sqrt{n} \rfloor$. Тогда в силу уже доказанного $S_{m^2(n)}/m^2(n) \toAS 0$ при $n \to \infty$. Обозначим

    $$ \zeta_n = \frac{S_n}{m^2(n)} - \frac{S_{m^2(n)}}{m^2(n)} $$
    и проверим, что $\zeta_n \toAS 0$. Для любого $\varepsilon > 0$
    $$ \P(|\zeta_n| \ge \varepsilon) \le |\text{неравенство Чебышева}| \le \frac{\D(S_n - S_{m^2(n)})}{\varepsilon^2 m^4(n)} =$$ 
    $$=|\text{некоррелированность}| = \frac{\sum_{k=m^2(n)}^n \D\xi_k}{\varepsilon^2 m^4(n)} \le \frac{C(n - m^2(n))}{\varepsilon^2 m^4(n)}. $$
    Далее, заметим, что $m^4(n) \sim n^2$ и $n - m^2(n) \le n - (\sqrt{n} - 1)^2 \le 2\sqrt{n}$. Стало быть, $\P(|\zeta_n| \ge \varepsilon) = O(n^{-3/2})$. Следовательно, ряд $\sum_{n=1}^\infty \P(|\zeta_n| \ge \varepsilon)$ сходится и $\zeta_n \toAS 0$.

    Подведем итоги. Мы показали, что $S_{m^2(n)}/m^2(n) \toAS 0$ и $\zeta_n \toAS 0$. Значит, $S_n/m^2(n) \toAS 0$. Осталось заметить, что $m^2(n) \sim n$.
}
